% !TeX root = ../../Book3.tex
\section{覆盖映射}
\label{b3:sec:1101}

多值函数可以放到一张覆盖空间上变成单值。覆盖的关键条件不是仅有局部逆，而是目标的每个小邻域都被完整、互不相交的各层覆盖。这个条件保证路径与同伦可以一直提升。

\subsection{覆盖映射}
\label{b3:subsec:110101}

\begin{definition}\label{b3:def:11-covering}
连通空间之间的连续满射 \(p:X\rightarrow Y\) 称为\term[fugai-yingshe]{覆盖映射}[covering map]，若每个 \(y\in Y\) 有开邻域 \(V\)，使 \(p^{-1}(V)\) 是互不相交开集 \(V_\alpha\) 的并，且每个 \(p|_{V_\alpha}:V_\alpha\rightarrow V\) 为同胚。这些 \(V_\alpha\) 称为层，\(V\) 称为均匀覆盖邻域。
\end{definition}
本章空间均为 Hausdorff、第二可数的曲面或区域，因此局部路径连通，路径提升所需的局部条件成立。纤维 \(p^{-1}(y)\) 离散，其点数在连通底空间上固定，允许无穷。这个点数称为层数；有限层数不意味着覆盖空间紧，底空间还须紧才有这一推论。

\ProseParagraphBreak
一个局部同胚未必是覆盖，因为某层可能只覆盖目标小邻域的一部分。例如 \(e^z:\mathbb C\setminus\{0\}\rightarrow\mathbb C^*\) 满射且导数从不为零，但1附近有一个被删掉中心的局部对数层，无法完整覆盖该邻域。因此它不是覆盖。完整指数 \(e^z:\mathbb C\rightarrow\mathbb C^*\) 才是覆盖。

\ProseParagraphBreak
局部同胚与覆盖的差别可以用层是否完整来判断。每一点有局部逆只说明一层在该点存在；覆盖要求同一个目标小邻域同时容纳所有原像层，并且没有层在邻域内突然结束或逃逸。开区间 \((0,3/2)\) 上的 \(t\mapsto e^{2\pi\mathrm{i}t}\) 是到圆周的满射局部同胚，但不是覆盖：在某些目标点有两层，在某些点只有一层。有限端点截断了一个本来应完整的层。这个实一维例子足以显示定义中的全部原像条件不是冗余的。

\ProseParagraphBreak
覆盖的层数在连通目标上恒定，允许无限。沿任一连接两点的路径提升，可把起始纤维的每一点送到终点纤维；逆路径给逆映射，因而纤维基数相同。这个论证不要求有限层，也不要求目标单连通。不同路径可能给不同纤维双射，闭路造成的置换正是覆盖的 monodromy 资料。

\subsection{路径提升}
\label{b3:subsec:110102}

\begin{theorem}[路径提升]\label{b3:thm:11-path-lifting}
给定覆盖 \(p:X\rightarrow Y\)、连续路径 \(\gamma\) 及 \(x_0\in p^{-1}(\gamma(0))\)，存在唯一连续路径 \(\widetilde\gamma\) 满足 \(p\circ\widetilde\gamma=\gamma\)、\(\widetilde\gamma(0)=x_0\)。
\end{theorem}
\begin{proof}
用均匀覆盖邻域覆盖紧路径像，按 Lebesgue 数把参数分成有限小段，每段像落在一个邻域中。首段使用包含 \(x_0\) 的层的逆映射；它的终点确定下一段所用层，依次得到连续提升。两提升若在一个位置相同，则在附近使用同一层逆而相同；相同位置集合开且闭，且包含0，故由区间连通性处处相同。
\end{proof}
路径每绕一个孔，提升终点可移到纤维的另一个点。起点选择是必要数据；若只给底路径，有限层覆盖通常有多条不同提升。闭路的提升也未必闭合，指数覆盖的单位圆提升从0走到 \(2\pi\mathrm{i}\)。

\subsection{同伦提升}
\label{b3:subsec:110103}

\begin{theorem}[同伦提升]\label{b3:thm:11-homotopy-lifting}
若 \(H:[0,1]^2\rightarrow Y\) 为路径同伦，给定 \(H(0,\cdot)\) 的一个提升，便有唯一提升 \(\widetilde H\) 延伸它。端点固定同伦的提升在给定起点后具有相同终点。
\end{theorem}
\begin{proof}
将参数方形分为有限小格，使每格像落在均匀覆盖邻域。已知第一侧提升后，逐格用包含一个已知顶点的层逆定义格上提升。公共边的两个定义是同一底路径的提升，且在端点相同，路径提升唯一性使它们一致，故所有格拼成连续提升。唯一性同样由局部层逆和连通性得到。若端点固定，提升终点沿 \(s\) 形成落在一个离散纤维中的连续路径，因此恒定。
\end{proof}
这使底空间基本群作用在纤维上：沿闭路提升取末点。作用只依赖同伦类，但它未必与甲板群相同；甲板群是整张覆盖的对称，而纤维作用是路径传递。

\begin{theorem}[映射提升的判别]\label{b3:thm:11-map-lifting}
设 \(p:X\to Y\) 是连通曲面之间的覆盖，\(F:Z\to Y\) 连续，\(Z\) 是连通曲面。选 \(z_0\in Z\)、\(x_0\in X\)，使 \(p(x_0)=F(z_0)=y_0\)。存在满足 \(\widetilde F(z_0)=x_0\) 的提升，当且仅当
\[
 F_*\pi_1(Z,z_0)\subset p_*\pi_1(X,x_0).
\]
满足基点条件的提升唯一；若各曲面带复结构，\(p,F\) 全纯，则提升全纯。
\end{theorem}
\begin{proof}
若有提升，闭路先在 \(Z\) 映到 \(X\) 再投到 \(Y\)，故群包含必要。反向对 \(z\) 选从 \(z_0\) 到 \(z\) 的路径 \(\gamma\)，把 \(F\circ\gamma\) 从 \(x_0\) 提升，以终点定义 \(\widetilde F(z)\)。两条路径拼成闭路，其像类属于 \(p_*\pi_1(X)\)。因此它与某条 \(X\) 闭路的投影同伦，利用同伦提升，像闭路的提升也闭合。两条原路径的提升终点相同，定义不依赖选择。每个 \(z\) 附近选小连通邻域，使其像落在覆盖均匀邻域；这时提升等于固定局部逆与 \(F\) 的复合，故连续。两提升相等的点集在局部逆图内开，其不相等的点集也开，连通性给唯一。全纯情形的同一局部复合表明提升全纯。
\end{proof}
特别地，单连通源的群为0，所以总可提升。第十三章把圆盘上的全纯映射提升到目标曲面的圆盘覆盖时，用的正是这个特例。源仅为一个圆环时没有这样的自动结论：映射在基本圈上绕行的次数必须属于目标覆盖对应的子群。

\ProseParagraphBreak
例如 \(p(w)=w^m:\mathbb C^*\to\mathbb C^*\)，其基本群像是 \(m\mathbb Z\)。一个无零全纯函数 \(F:U\to\mathbb C^*\) 提升到这张覆盖，恰是找到 \(h^m=F\)；判别条件就是所有像闭路的绕数被 \(m\) 整除。这与第十章根的周期条件相同，两个观点分别把它写成群包含和对数导数积分。

\subsection{甲板变换}
\label{b3:subsec:110104}

\begin{definition}\label{b3:def:11-deck}
满足 \(p\circ T=p\) 的覆盖空间自同胚称为\term[jiaban-bianhuan]{甲板变换}[deck transformation]，全部甲板变换组成群 \(\operatorname{Deck}(p)\)。
\end{definition}
\begin{proposition}\label{b3:prop:11-deck-free}
连通覆盖的甲板变换若固定一点，便是恒等。每点有小邻域与其任何非恒等甲板像不交。
\end{proposition}
\begin{proof}
连接固定点到任意点的路径，经 \(T\) 变换后仍是同一底路径的提升且起点相同，唯一性使末点不变，所以 \(T\) 恒等。对均匀覆盖邻域中的一层 \(W\)，若 \(T(W)\) 与 \(W\) 相交，它们同为该邻域的层，故相同；层上投影单射使 \(T\) 在 \(W\) 恒等，再由第一结论为全域恒等。
\end{proof}
任意覆盖的甲板群不必在纤维上传递。若它确实传递，称该覆盖为正规覆盖；万有覆盖将满足这一性质。不要把“局部各层相同”直接解释为存在交换任意两层的全局甲板变换。
